\section{Method audit and native objective}
\label{app:methods}

This appendix specifies the historical native objective inspected at the paper's
repository snapshot. It is not the specification of a completed InnovationZero
experiment. The principal implementation is
\texttt{deploy/dl-flywheel/pipelines/continuous/}\allowbreak
\texttt{neural\_native\_batch.py}; model and sequence-loss primitives live in its
\texttt{consumer/dltrain/} directory. The macro-batch step is in
\texttt{neural\_native\_train.py}; the subsequent anchor is in
\texttt{pipelines/continuous\_v6/anchored\_step.py}. The repository report and
claim ledger give the retained run evidence~\citep{clover,jope}.

\subsection{Quality, support, and reduction}

The target scale is $s_c=1.4826\,\mathrm{MAD}(y)$ when positive, otherwise
$\mathrm{IQR}(y)/1.349$ when positive, otherwise the population standard
deviation. Set
\begin{align}
 \varepsilon_c&=\max\{10^{-12}\max(1,|\bar y|),10^{-6}s_c\},\nonumber\\
 q_{\mathrm{opt}}&=\operatorname{clip}(q,-\log1000,\log1000),\nonumber\\
 q_{\mathrm{grid}}&=\operatorname{round}_{\mathrm{half-even}}(10^6q).
\end{align}
Here $q$ is the main paper's two-fold inner quality score. Missing baseline evidence remains
unknown. A completed-row path without fold-level detail can use its recorded
cross-validation log ratio under the native label contract; not every retained
row contains independently inspectable fold errors. Label provenance must travel
with the target.

For eligible rows $S_c^k$ of loss term $k$, define the reduction
\begin{equation}
 \mathcal{M}_k(v)=\frac{1}{|C_k|}\sum_{c\in C_k}
 \frac{1}{|S_c^k|}\sum_{i\in S_c^k}v_i,
 \qquad C_k=\{c:|S_c^k|>0\}.
\end{equation}
An empty support has zero numerator and denominator and a graph-connected zero
loss. It does not create a pseudo-observation. Microbatch accumulation divides
each term by the \emph{macro-batch} supported-context total, preserving the
declared weighting when sequence lengths change the packing. Consumed targets,
head outputs, and gradients must be finite.

For a sequence of length $T_i$ including boundary symbols, native CE is
\begin{equation}
 L_{\mathrm{CE}}=\mathcal{M}_{\mathrm{CE}}\!\left(
 -\frac{1}{T_i-1}\sum_{t=1}^{T_i-1}
 \log p_\theta(a_{it}\mid c,a_{i,<t})\right).
\end{equation}
The beginning symbol is not a prediction target; the end symbol is. The
conditional distribution uses the training legality mask. Actual CE-eligible
sequences receive uniform weight within their context. This differs from the
teacher-softmax weighting in the separate generic grouped consumer. The quality
term is
$L_q=\mathcal{M}_q(\operatorname{Huber}_1(\hat q_i-q_{\mathrm{opt},i}))$.

\subsection{Ranking, feasibility, and censored runtime}

Let $r_i$ be ascending within-context midrank percentile, with ties averaged and
singleton rank $1/2$. Eligible ordered pairs have
$q_{\mathrm{grid},i}>q_{\mathrm{grid},k}$. Their weights are
\begin{equation}
 w_{ik}=\min(|q_{\mathrm{opt},i}-q_{\mathrm{opt},k}|,1)
 \bigl(1+\mathbf{1}\{\max(r_i,r_k)\geq0.9\}\bigr).
\end{equation}
The Bradley--Terry term averages
$\sum_{i,k}w_{ik}\operatorname{softplus}(\hat q_k-\hat q_i)/\sum_{i,k}w_{ik}$
over contexts with positive total weight. Pairwise eligibility requires measured
quality; tied grid values do not form ordered pairs. The list term applies only
to contexts with at least three pairwise-eligible programs:
\begin{equation}
 L_{\mathrm{list}}=\operatorname{mean}_c
 \left[-\sum_i \operatorname{softmax}_i(r_i/0.25)
 \log\operatorname{softmax}_i(\hat q_i)\right].
\end{equation}
Feasibility is three-class cross-entropy over recorded completed, failed, or
timeout labels, reduced by $\mathcal{M}_{\mathrm{feas}}$. Unknown outcomes are
not relabeled as failures. For uncensored supported runtime $t_i$, the loss is
$\operatorname{Huber}_1(\hat t_i-\log(1+t_i))$. For right-censored timeout
observations it is $[\operatorname{ReLU}(\log(1+t_{\mathrm{bound}})-\hat t_i)]^2$;
the historical 90-second bound gives $\log91$. The runtime term uses
$\mathcal{M}_{\mathrm{rt}}$. These six terms combine as in
the main paper's six-term $L_6$ equation.

\subsection{Sampling, anchoring, and interpretation}

The recorded batching contract takes up to 16 training contexts, at most one per
parent lineage, with an eligible teacher and a rotating window of up to eight
measured programs per context. It packs at most 8,192 padded tokens without
truncating complete programs. Reported training data include 7,120 complete
contexts: 6,857 training and 263 validation, associated with 1,014 training and
54 validation parent lineages. These are historical dataset roles, not the
independent family count of a new confirmatory campaign.

The v2 objective adds
\begin{equation}
 L_{\mathrm{v2}}=L_6+
 0.2\operatorname{mean}_{c\in K}
 D_{\mathrm{KL}}(\pi_\theta(\cdot\mid c)\Vert\pi_{\mathrm{v1}}(\cdot\mid c)).
\end{equation}
The anchor is exact over finite deployed support, with temperature one and
dropout off. The inspected configuration selects the batch's first context for
this term; an exact conditional KL is not an all-context population estimate.
Reported \texttt{loss} remains $L_6$; KL and the total objective are separate
fields. Training likelihood and deployed decoding use different normalizations,
which is another reason that lower CE need not improve deployed utility.

An expanded-support successor cannot retain this anchor unchanged. If
$\pi_\theta(a\mid c)>0$ for an action with
$\pi_{\mathrm{v1}}(a\mid c)=0$, then
$D_{\mathrm{KL}}(\pi_\theta\Vert\pi_{\mathrm{v1}})=\infty$.
An open compositional decoder therefore needs a reference with compatible
support, or an explicitly normalized mixture assigning positive probability
to every permitted action, and a separately declared regularization objective.
If several decoding trajectories normalize to the same program, its probability
is the sum over those trajectories. A tractable trajectory-level KL must be
identified as such; it is not automatically the KL between canonical-program
distributions. These are requirements for a proposed successor, not claims that
the existing finite-catalog implementation violates its own support contract.

The 2,048-value Candle ABI is documented in
\nolinkurl{crates/jope-policy/src/context.rs} and assembled by
\texttt{context\_producer.rs}; its inference architecture is in
\texttt{model.rs}. Hyperion's 512-dimensional signature is implemented in
\nolinkurl{crates/hyperion/src/hyperion_v7/model.rs}.
\nolinkurl{crates/jope-pipeline/src/neural_search.rs} describes policy seeding
followed by bounded search and explicitly qualifies arithmetic parity.
These distinct sources cannot be combined into a fictional single checkpoint.
The future \texttt{solution\_embedding} needs its own dimension contract,
checkpoint identity, decomposition training corpus, decoder support, and
controlled ablations before it can explain an observed result.

\subsection{Proposed training comparison}

A deployment-aligned arm can estimate the main paper's first-proposal objective $J_1$
from on-policy training-family outcomes using
\begin{equation}
 \widehat{\nabla_\theta J_1}=
 \frac{1}{m}\sum_{i=1}^{m}
 (u_i-v(c_i))\nabla_\theta\log\pi_\theta(a_i\mid c_i),
\end{equation}
where the baseline $v(c_i)$ does not depend on the sampled action, and gradients
do not pass through measured utility or the baseline. The estimator is a
proposed experiment, not the current native objective. Historical offline
records are not on-policy samples; using them needs an explicitly justified
off-policy method rather than substituting them into this expression.

For the proposed solution embedding, first decompose training-only procedures
into versioned structural objects, then learn reconstruction and measured-behavior
representations, and finally compare conditional proposal training against
direct conditioning without the bottleneck. Reconstruction success alone does
not establish useful synthesis. Freeze the decomposition version and source
archive before confirmatory families are exposed; compare equal-compute controls,
shuffled contexts, nearest-program retrieval, catalog selection, and removal of
the decomposition stage. These interventions distinguish a useful learned
representation from extra capacity or a larger lookup table.
